\chapter{Arguments}

\section{Why This Matters}
When you pass data into a function, Python needs to know \textit{which} parameter gets \textit{which} piece of data. Understanding Positional, Keyword, and Default arguments allows you to write incredibly flexible ML functions.

\section{Positional vs. Keyword Arguments}
By default, arguments are mapped to parameters based on their \textbf{position} (order).
Instead of relying on order, you can explicitly state which parameter gets which value by using the parameter's name. This is called a \textbf{keyword argument}.

\begin{lstlisting}
def describe_model(name, layers):
    print(f"Model: {name}, Layers: {layers}")

# Positional (Order matters)
describe_model("ResNet", 50) 

# Keyword (Order doesn't matter)
describe_model(layers=50, name="ResNet") 
\end{lstlisting}
\textit{Rule: If you mix positional and keyword arguments, positional arguments \textbf{MUST} come first!}

\section{Default Arguments}
You can provide a \textbf{default value} in the function definition. If the user doesn't provide an argument for it, Python uses the default.
\begin{lstlisting}
def train_model(model_name, epochs=10):
    print(f"Training {model_name} for {epochs} epochs.")

train_model("YOLOv8")              # Uses default (10)
train_model("YOLOv8", epochs=100)  # Overrides default (100)
\end{lstlisting}
\textit{Rule: Parameters with default values MUST come \textbf{after} parameters without default values.}

\section{\texttt{*args} and \texttt{**kwargs} (The Advanced Magic)}
Sometimes you don't know how many arguments the user will pass. 
\begin{itemize}
    \item \texttt{*args}: Allows you to pass any number of \textbf{positional} arguments. They are packed into a \textbf{Tuple}.
    \item \texttt{**kwargs}: Allows you to pass any number of \textbf{keyword} arguments. They are packed into a \textbf{Dictionary}.
\end{itemize}

\begin{lstlisting}
def calculate_total_loss(*losses):
    return sum(losses)
print(calculate_total_loss(0.5, 0.2, 0.1)) # 0.8
\end{lstlisting}

\mlconnection{When you use a library like \texttt{scikit-learn} or \texttt{TensorFlow}, functions have dozens of parameters. They rely heavily on \textbf{Default Arguments} (so you don't have to specify all 20 settings every time) and \textbf{Keyword Arguments} (so you know exactly what you are modifying). Wrappers heavily use \texttt{**kwargs} to pass arbitrary parameters down to lower-level functions.}

\section{Solutions to Exercises}

\textbf{Exercise 1: The Configurator}
\begin{lstlisting}
def setup_server(host, port=8080, debug=False):
    print(f"Server starting on {host}:{port} | Debug mode: {debug}")

setup_server("localhost")
setup_server("192.168.1.1", 9000)
setup_server("production.com", debug=True)
\end{lstlisting}

\vspace{1em}
\textbf{Exercise 2: Keyword Enforcement}
\begin{lstlisting}
def train_network(learning_rate, batch_size, epochs, momentum):
    print(f"LR: {learning_rate}, Batch: {batch_size}, Epochs: {epochs}, Momentum: {momentum}")

train_network(learning_rate=0.001, batch_size=64, epochs=100, momentum=0.9)
\end{lstlisting}

\vspace{1em}
\textbf{Exercise 3: The \texttt{*args} Aggregator}
\begin{lstlisting}
def average_scores(*args):
    return sum(args) / len(args)

print(average_scores(85, 90, 95, 100)) # 92.5
\end{lstlisting}
